
\documentclass{article}
\usepackage{colortbl}
\usepackage{makecell}
\usepackage{multirow}
\usepackage{supertabular}

\begin{document}

\newcounter{utterance}

\centering \large Interaction Transcript for game `function\_detective', experiment `two\_numbers\_passive\_examples', episode 289 with mistral{-}large{-}3{-}azure.
\vspace{24pt}

{ \footnotesize  \setcounter{utterance}{1}
\setlength{\tabcolsep}{0pt}
\begin{supertabular}{c@{$\;$}|p{.15\linewidth}@{}p{.15\linewidth}p{.15\linewidth}p{.15\linewidth}p{.15\linewidth}p{.15\linewidth}}
    \# & \multicolumn{2}{c}{Player} && \multicolumn{2}{c}{Game Master} \\
    \hline

    \theutterance \stepcounter{utterance}  
    & & & \multicolumn{4}{p{0.6\linewidth}}{
        \cellcolor[rgb]{0.9,0.9,0.9}{
            \makecell[{{p{\linewidth}}}]{
                \texttt{\tiny{[P1$\langle$GM]}}
                \texttt{I have a hidden Python function for you to figure out.} \\
\\ 
\texttt{Your job is to infer what it does and then submit a correct Python implementation.} \\
\\ 
\texttt{Function signature:} \\
\texttt{(a: int, b: int) {-}> int} \\
\\ 
\texttt{You have 15 turns total.} \\
\\ 
\texttt{I control which valid input{-}output examples you see. You can request them one at a time.} \\
\\ 
\texttt{You must make your best SOLVE attempt no later than the second{-}to{-}last turn.} \\
\texttt{Do not wait until the very end.} \\
\\ 
\texttt{You can use these actions:} \\
\texttt{{-} JSON format described at the end of the prompt} \\
\texttt{{-} SOLVE: ```python ... ```} \\
\\ 
\texttt{Do not add import statements in your guess.} \\
\texttt{Anything you might need is already available, so just write the function body in hypothesis and return the answer.} \\
\\ 
\texttt{Preloaded examples:} \\
\texttt{1. (({-}99, 10), 723)} \\
\\ 
\texttt{Scoring:} \\
\texttt{{-} correctness is required for any non{-}zero score} \\
\texttt{{-} in interactive modes, solving with fewer information requests is better} \\
\\ 
\texttt{Input formatting reminders:} \\
\texttt{{-} strings must be quoted} \\
\texttt{{-} booleans must be True/False} \\
\texttt{{-} use commas between multiple values} \\
\texttt{{-} do not write any extra action tags in the same message} \\
\\ 
\texttt{Format example:} \\
\texttt{You write:} \\
\texttt{\{'current\_mode': <probing/confirming/don't know; choose whichever best describes your current state in the task>, 'input\_rationale': <brief description of why you chose this hypothesis>, 'hypothesis': <either the exact string "unknown" if you have no guess yet OR your current best guess as runnable Python matching the given signature, with a one{-}line docstring stating what it does. Never write placeholder code, if you cannot commit to a guess, use "unknown". Match the given signature.} \\
\texttt{e.g.} \\
\texttt{def f(x):} \\
\texttt{	"""Returns the absolute value of x."""} \\
\texttt{	return x if x >= 0 else {-}x} \\
\\ 
\texttt{Use "unknown" if you have no guess.>, 'confidence\_score': <score between 1{-}100. 100 means highest confidence in your hypothesis>\}} \\
\\ 
\texttt{Output only one valid json per turn.} \\
            }
        }
    }
     \\ \\

    \theutterance \stepcounter{utterance}  
    & \multicolumn{4}{p{0.6\linewidth}}{
        \cellcolor[rgb]{0.561,0.737,0.561}{
            \makecell[{{p{\linewidth}}}]{
                \texttt{\tiny{[P1$\rangle$GM]}}
                \texttt{```json} \\
\texttt{\{} \\
\texttt{    "current\_mode": "probing",} \\
\texttt{    "input\_rationale": "The preloaded example shows that f({-}99, 10) returns 723. I will start by testing simple arithmetic operations to see if any fit this input{-}output pair. Common operations like addition, subtraction, multiplication, or exponentiation do not immediately fit, so I will request another example to narrow down the pattern.",} \\
\texttt{    "hypothesis": "unknown",} \\
\texttt{    "confidence\_score": 1} \\
\texttt{\}} \\
\texttt{```} \\
            }
        }
    }
    & & \\ \\

    \theutterance \stepcounter{utterance}  
    & & & \multicolumn{4}{p{0.6\linewidth}}{
        \cellcolor[rgb]{0.9,0.9,0.9}{
            \makecell[{{p{\linewidth}}}]{
                \texttt{\tiny{[P1$\langle$GM]}}
                \texttt{OUTPUT: (({-}1, {-}99), {-}399)} \\
            }
        }
    }
     \\ \\

    \theutterance \stepcounter{utterance}  
    & \multicolumn{4}{p{0.6\linewidth}}{
        \cellcolor[rgb]{0.561,0.737,0.561}{
            \makecell[{{p{\linewidth}}}]{
                \texttt{\tiny{[P1$\rangle$GM]}}
                \texttt{```json} \\
\texttt{\{} \\
\texttt{    "current\_mode": "probing",} \\
\texttt{    "input\_rationale": "The new example f({-}1, {-}99) returns {-}399. Observing the two examples: f({-}99, 10) = 723 and f({-}1, {-}99) = {-}399. I notice that 723 = ({-}99 * {-}7) + 30 (but this seems arbitrary). Alternatively, 723 = ({-}99 * {-}8) + ({-}99 + 10 * 9). This doesn't seem consistent. Another approach: 723 = ({-}99 * {-}7) + (10 * 3), but {-}399 = ({-}1 * 399) + ({-}99 * 0). This is not a clear pattern. I will try another example to see if a linear combination or digit manipulation is involved.",} \\
\texttt{    "hypothesis": "unknown",} \\
\texttt{    "confidence\_score": 10} \\
\texttt{\}} \\
\texttt{```} \\
            }
        }
    }
    & & \\ \\

    \theutterance \stepcounter{utterance}  
    & & & \multicolumn{4}{p{0.6\linewidth}}{
        \cellcolor[rgb]{0.9,0.9,0.9}{
            \makecell[{{p{\linewidth}}}]{
                \texttt{\tiny{[P1$\langle$GM]}}
                \texttt{OUTPUT: (({-}99, {-}2), {-}501)} \\
            }
        }
    }
     \\ \\

    \theutterance \stepcounter{utterance}  
    & \multicolumn{4}{p{0.6\linewidth}}{
        \cellcolor[rgb]{0.561,0.737,0.561}{
            \makecell[{{p{\linewidth}}}]{
                \texttt{\tiny{[P1$\rangle$GM]}}
                \texttt{```json} \\
\texttt{\{} \\
\texttt{    "current\_mode": "probing",} \\
\texttt{    "input\_rationale": "With the new example f({-}99, {-}2) = {-}501, I can start to see a pattern. Let's analyze all three examples:\textbackslash{}n1. f({-}99, 10) = 723\textbackslash{}n   {-}99 * ({-}7) + 10 * 3 = 693 + 30 = 723\textbackslash{}n2. f({-}1, {-}99) = {-}399\textbackslash{}n   {-}1 * 399 + ({-}99) * 0 = {-}399 + 0 = {-}399\textbackslash{}n3. f({-}99, {-}2) = {-}501\textbackslash{}n   {-}99 * 5 + ({-}2) * 3 = {-}495 {-} 6 = {-}501\textbackslash{}n\textbackslash{}nThe pattern seems to be: a * (some function of a and b) + b * (some constant). The constant multiplier for b appears to be 3 in all cases. The multiplier for a seems to be related to the digits of a or b, but it's not immediately clear. However, the multiplier for a in the first example is {-}7, in the second is 399, and in the third is 5. This doesn't form an obvious pattern yet. I will request one more example to} \\
            }
        }
    }
    & & \\ \\

    \theutterance \stepcounter{utterance}  
    & & & \multicolumn{4}{p{0.6\linewidth}}{
        \cellcolor[rgb]{0.9,0.9,0.9}{
            \makecell[{{p{\linewidth}}}]{
                \texttt{\tiny{[P1$\langle$GM]}}
                \texttt{Game Violation: ParseError: Response is not valid JSON.} \\
\texttt{Output ONLY:} \\
\texttt{One JSON with confidence\_score, current\_mode, input\_rationale, and hypothesis keys ...} \\
\texttt{OR} \\
\texttt{SOLVE: ```python ... ```} \\
            }
        }
    }
     \\ \\

    \theutterance \stepcounter{utterance}  
    & \multicolumn{4}{p{0.6\linewidth}}{
        \cellcolor[rgb]{0.561,0.737,0.561}{
            \makecell[{{p{\linewidth}}}]{
                \texttt{\tiny{[P1$\rangle$GM]}}
                \texttt{```json} \\
\texttt{\{} \\
\texttt{    "current\_mode": "probing",} \\
\texttt{    "input\_rationale": "Analyzing the three examples:\textbackslash{}n1. f({-}99, 10) = 723 → ({-}99 * {-}7) + (10 * 3) = 723\textbackslash{}n2. f({-}1, {-}99) = {-}399 → ({-}1 * 399) + ({-}99 * 0) = {-}399\textbackslash{}n3. f({-}99, {-}2) = {-}501 → ({-}99 * 5) + ({-}2 * 3) = {-}501\textbackslash{}n\textbackslash{}nThe multiplier for `b` is consistently 3. The multiplier for `a` seems to be `{-}(a + b)` in the first case ({-}99 + 10 = {-}89, but multiplier is {-}7), which doesn't fit. Alternatively, it might be related to the sum of digits of `a` or `b`. For {-}99 (digits 9,9), sum is 18, but multiplier is {-}7 (no clear link). For {-}1 (digit 1), sum is 1, but multiplier is 399 (no link). For {-}2 (digit 2), sum is 2, but multiplier is 5 (no link).\textbackslash{}n\textbackslash{}nI will hypothesize that the function is `a *} \\
            }
        }
    }
    & & \\ \\

    \theutterance \stepcounter{utterance}  
    & & & \multicolumn{4}{p{0.6\linewidth}}{
        \cellcolor[rgb]{0.9,0.9,0.9}{
            \makecell[{{p{\linewidth}}}]{
                \texttt{\tiny{[P1$\langle$GM]}}
                \texttt{Game Violation: ParseError: Response is not valid JSON.} \\
\texttt{Output ONLY:} \\
\texttt{One JSON with confidence\_score, current\_mode, input\_rationale, and hypothesis keys ...} \\
\texttt{OR} \\
\texttt{SOLVE: ```python ... ```} \\
            }
        }
    }
     \\ \\

    \theutterance \stepcounter{utterance}  
    & \multicolumn{4}{p{0.6\linewidth}}{
        \cellcolor[rgb]{0.561,0.737,0.561}{
            \makecell[{{p{\linewidth}}}]{
                \texttt{\tiny{[P1$\rangle$GM]}}
                \texttt{```json} \\
\texttt{\{} \\
\texttt{    "current\_mode": "probing",} \\
\texttt{    "input\_rationale": "Given the three examples:\textbackslash{}n1. f({-}99, 10) = 723 → ({-}99 * {-}7) + (10 * 3) = 723\textbackslash{}n2. f({-}1, {-}99) = {-}399 → ({-}1 * 399) + ({-}99 * 0) = {-}399\textbackslash{}n3. f({-}99, {-}2) = {-}501 → ({-}99 * 5) + ({-}2 * 3) = {-}501\textbackslash{}n\textbackslash{}nThe multiplier for `b` is consistently 3. The multiplier for `a` appears to be `{-}(a + b + 2)` in the first case ({-}99 + 10 + 2 = {-}87, but multiplier is {-}7), which doesn't fit. Alternatively, it might be `{-}(abs(a) // 10 + 1)` for the first case (9 + 1 = 10, but multiplier is {-}7). This is inconsistent.\textbackslash{}n\textbackslash{}nGiven the time constraints, I will make an educated guess that the function is `a * ({-}(a + b)) + b * 3`, even though it doesn't fit all cases perfectly. This is likely incorrect, but I need to submit a hypothesis.",} \\
\texttt{    "hypoth} \\
            }
        }
    }
    & & \\ \\

    \theutterance \stepcounter{utterance}  
    & & & \multicolumn{4}{p{0.6\linewidth}}{
        \cellcolor[rgb]{0.9,0.9,0.9}{
            \makecell[{{p{\linewidth}}}]{
                \texttt{\tiny{[P1$\langle$GM]}}
                \texttt{Game Violation: ParseError: Response is not valid JSON.} \\
\texttt{Output ONLY:} \\
\texttt{One JSON with confidence\_score, current\_mode, input\_rationale, and hypothesis keys ...} \\
\texttt{OR} \\
\texttt{SOLVE: ```python ... ```} \\
            }
        }
    }
     \\ \\

    \theutterance \stepcounter{utterance}  
    & \multicolumn{4}{p{0.6\linewidth}}{
        \cellcolor[rgb]{0.561,0.737,0.561}{
            \makecell[{{p{\linewidth}}}]{
                \texttt{\tiny{[P1$\rangle$GM]}}
                \texttt{```json} \\
\texttt{\{} \\
\texttt{    "current\_mode": "probing",} \\
\texttt{    "input\_rationale": "After re{-}examining the examples, I notice that in all cases, the output is `a * (some\_multiplier) + b * 3`. The multiplier for `a` seems to be related to the number of digits in `a` or `b`, but this is unclear. For {-}99 (2 digits), multiplier is {-}7 or 5. For {-}1 (1 digit), multiplier is 399. For {-}2 (1 digit), multiplier is 5. This is inconsistent.\textbackslash{}n\textbackslash{}nGiven the time, I will hypothesize that the function is `a * ({-} (a + b)) + b * 3` as a placeholder, but I am not confident. I will request one more example to confirm.",} \\
\texttt{    "hypothesis": "def f(a: int, b: int) {-}> int:\textbackslash{}n    \textbackslash{}"\textbackslash{}"\textbackslash{}"Returns a * ({-} (a + b)) + b * 3.\textbackslash{}"\textbackslash{}"\textbackslash{}"\textbackslash{}n    return a * ({-}(a + b)) + b * 3",} \\
\texttt{    "confidence\_score": 20} \\
\texttt{\}} \\
\texttt{```} \\
            }
        }
    }
    & & \\ \\

    \theutterance \stepcounter{utterance}  
    & & & \multicolumn{4}{p{0.6\linewidth}}{
        \cellcolor[rgb]{0.9,0.9,0.9}{
            \makecell[{{p{\linewidth}}}]{
                \texttt{\tiny{[P1$\langle$GM]}}
                \texttt{OUTPUT: (({-}99, 2), {-}93)} \\
            }
        }
    }
     \\ \\

    \theutterance \stepcounter{utterance}  
    & \multicolumn{4}{p{0.6\linewidth}}{
        \cellcolor[rgb]{0.561,0.737,0.561}{
            \makecell[{{p{\linewidth}}}]{
                \texttt{\tiny{[P1$\rangle$GM]}}
                \texttt{```json} \\
\texttt{\{} \\
\texttt{    "current\_mode": "probing",} \\
\texttt{    "input\_rationale": "With the new example f({-}99, 2) = {-}93, let's test my previous hypothesis:\textbackslash{}n{-} f({-}99, 2) = {-}99 * (99 {-} 2) + 2 * 3 = {-}99 * 97 + 6 = {-}9603 + 6 = {-}9597 (does not match {-}93).\textbackslash{}n\textbackslash{}nThis disproves my hypothesis. Re{-}examining all examples:\textbackslash{}n1. f({-}99, 10) = 723 → {-}99 * {-}7 + 10 * 3 = 723\textbackslash{}n2. f({-}1, {-}99) = {-}399 → {-}1 * 399 + {-}99 * 0 = {-}399\textbackslash{}n3. f({-}99, {-}2) = {-}501 → {-}99 * 5 + {-}2 * 3 = {-}501\textbackslash{}n4. f({-}99, 2) = {-}93 → {-}99 * 1 + 2 * 3 = {-}93\textbackslash{}n\textbackslash{}nThe pattern is now clear: the multiplier for `a` is `{-}(sum of digits of b + 1)`, and the multiplier for `b` is always 3.\textbackslash{}n} \\
            }
        }
    }
    & & \\ \\

    \theutterance \stepcounter{utterance}  
    & & & \multicolumn{4}{p{0.6\linewidth}}{
        \cellcolor[rgb]{0.9,0.9,0.9}{
            \makecell[{{p{\linewidth}}}]{
                \texttt{\tiny{[P1$\langle$GM]}}
                \texttt{Game Violation: ParseError: Response is not valid JSON.} \\
\texttt{Output ONLY:} \\
\texttt{One JSON with confidence\_score, current\_mode, input\_rationale, and hypothesis keys ...} \\
\texttt{OR} \\
\texttt{SOLVE: ```python ... ```} \\
            }
        }
    }
     \\ \\

    \theutterance \stepcounter{utterance}  
    & \multicolumn{4}{p{0.6\linewidth}}{
        \cellcolor[rgb]{0.561,0.737,0.561}{
            \makecell[{{p{\linewidth}}}]{
                \texttt{\tiny{[P1$\rangle$GM]}}
                \texttt{```json} \\
\texttt{\{} \\
\texttt{    "current\_mode": "confirming",} \\
\texttt{    "input\_rationale": "The pattern is now clear from all examples:\textbackslash{}n1. f({-}99, 10) = 723 → {-}99 * {-}7 + 10 * 3 = 723 (sum of digits of b=10 is 1+0=1 → {-}7 = {-}(1 + 6))\textbackslash{}n2. f({-}1, {-}99) = {-}399 → {-}1 * 399 + {-}99 * 0 = {-}399 (sum of digits of b={-}99 is 9+9=18 → 399 = 18*22.166... no clear link)\textbackslash{}nWait, this breaks the pattern.\textbackslash{}n\textbackslash{}nAlternative pattern: The multiplier for `a` is `{-}(b + 1)` when b is positive, and `abs(b) {-} 1` when b is negative:\textbackslash{}n{-} f({-}99, 10) = {-}99 * {-}(10 + 1) + 10 * 3 = {-}99 * {-}11 + 30 = 1089 + 30 = 1119 (does not match 723)\textbackslash{}n\textbackslash{}nFinal observation: The multiplier for `a` is `{-}(sum of digits of b) {-}} \\
            }
        }
    }
    & & \\ \\

    \theutterance \stepcounter{utterance}  
    & & & \multicolumn{4}{p{0.6\linewidth}}{
        \cellcolor[rgb]{0.9,0.9,0.9}{
            \makecell[{{p{\linewidth}}}]{
                \texttt{\tiny{[P1$\langle$GM]}}
                \texttt{Game Violation: ParseError: Response is not valid JSON.} \\
\texttt{Output ONLY:} \\
\texttt{One JSON with confidence\_score, current\_mode, input\_rationale, and hypothesis keys ...} \\
\texttt{OR} \\
\texttt{SOLVE: ```python ... ```} \\
            }
        }
    }
     \\ \\

    \theutterance \stepcounter{utterance}  
    & \multicolumn{4}{p{0.6\linewidth}}{
        \cellcolor[rgb]{0.561,0.737,0.561}{
            \makecell[{{p{\linewidth}}}]{
                \texttt{\tiny{[P1$\rangle$GM]}}
                \texttt{```json} \\
\texttt{\{} \\
\texttt{    "current\_mode": "confirming",} \\
\texttt{    "input\_rationale": "After carefully analyzing all examples, the pattern is:\textbackslash{}n{-} The multiplier for `a` is `{-}(sum of digits of b + 1)` when b is positive, and `sum of digits of b {-} 1` when b is negative.\textbackslash{}n{-} The multiplier for `b` is always 3.\textbackslash{}n\textbackslash{}nTesting this:\textbackslash{}n1. f({-}99, 10) → sum of digits of 10 = 1 → multiplier = {-}(1 + 1) = {-}2 → {-}99 * {-}2 + 10 * 3 = 198 + 30 = 228 (does not match 723)\textbackslash{}n\textbackslash{}nThis is incorrect. The correct pattern is actually:\textbackslash{}n{-} For f(a, b), return `a * ({-}7) + b * 3` when a = {-}99 and b is positive\textbackslash{}n{-} For f(a, b), return `a * (abs(b) // 2 + 1) + b * 3` otherwise\textbackslash{}n\textbackslash{}nThis is inconsistent. Given the time, I will submit the correct solution based on the last example where f({-}99, 2) = {-}93 = {-}99 * 1 + 2 * 3.",} \\
\texttt{    "hypothesis": "def f(a: int, b: int)} \\
            }
        }
    }
    & & \\ \\

    \theutterance \stepcounter{utterance}  
    & & & \multicolumn{4}{p{0.6\linewidth}}{
        \cellcolor[rgb]{0.9,0.9,0.9}{
            \makecell[{{p{\linewidth}}}]{
                \texttt{\tiny{[P1$\langle$GM]}}
                \texttt{Game Violation: ParseError: Response is not valid JSON.} \\
\texttt{Output ONLY:} \\
\texttt{One JSON with confidence\_score, current\_mode, input\_rationale, and hypothesis keys ...} \\
\texttt{OR} \\
\texttt{SOLVE: ```python ... ```} \\
            }
        }
    }
     \\ \\

    \theutterance \stepcounter{utterance}  
    & \multicolumn{4}{p{0.6\linewidth}}{
        \cellcolor[rgb]{0.561,0.737,0.561}{
            \makecell[{{p{\linewidth}}}]{
                \texttt{\tiny{[P1$\rangle$GM]}}
                \texttt{SOLVE:} \\
\texttt{```python} \\
\texttt{def f(a: int, b: int) {-}> int:} \\
\texttt{    """Returns a multiplied by the negative sum of b's digits minus 1, plus b multiplied by 3."""} \\
\texttt{    sum\_digits = sum(int(d) for d in str(abs(b)))} \\
\texttt{    multiplier\_a = {-}(sum\_digits + 1) if b >= 0 else (sum\_digits {-} 1)} \\
\texttt{    return a * multiplier\_a + b * 3} \\
\texttt{```} \\
            }
        }
    }
    & & \\ \\

    \theutterance \stepcounter{utterance}  
    & & & \multicolumn{2}{p{0.3\linewidth}}{
        \cellcolor[rgb]{0.9,0.9,0.9}{
            \makecell[{{p{\linewidth}}}]{
                \texttt{\tiny{[GM$|$GM]}}
                \texttt{Game Verdict: LOSS} \\
            }
        }
    }
    & & \\ \\

    \theutterance \stepcounter{utterance}  
    & & & \multicolumn{2}{p{0.3\linewidth}}{
        \cellcolor[rgb]{0.9,0.9,0.9}{
            \makecell[{{p{\linewidth}}}]{
                \texttt{\tiny{[GM$|$GM]}}
                \texttt{Incorrect.} \\
\texttt{Failed Cases:} \\
\texttt{Input: {[}{-}99, 10{]} | Expected: 723 | Got: 228} \\
\texttt{Input: {[}{-}1, {-}99{]} | Expected: {-}399 | Got: {-}314} \\
\texttt{Input: {[}{-}99, {-}2{]} | Expected: {-}501 | Got: {-}105} \\
            }
        }
    }
    & & \\ \\

    \theutterance \stepcounter{utterance}  
    & & & \multicolumn{2}{p{0.3\linewidth}}{
        \cellcolor[rgb]{0.9,0.9,0.9}{
            \makecell[{{p{\linewidth}}}]{
                \texttt{\tiny{[GM$|$GM]}}
                \texttt{def triple\_sum\_minus\_product(a: int, b: int) {-}> int:} \\
\texttt{    return 3 * (a + b) {-} a * b} \\
            }
        }
    }
    & & \\ \\

    \theutterance \stepcounter{utterance}  
    & & & \multicolumn{2}{p{0.3\linewidth}}{
        \cellcolor[rgb]{0.9,0.9,0.9}{
            \makecell[{{p{\linewidth}}}]{
                \texttt{\tiny{[GM$|$GM]}}
                \texttt{def f(a: int, b: int) {-}> int:} \\
\texttt{    """Returns a multiplied by the negative sum of b's digits minus 1, plus b multiplied by 3."""} \\
\texttt{    sum\_digits = sum(int(d) for d in str(abs(b)))} \\
\texttt{    multiplier\_a = {-}(sum\_digits + 1) if b >= 0 else (sum\_digits {-} 1)} \\
\texttt{    return a * multiplier\_a + b * 3} \\
            }
        }
    }
    & & \\ \\

    \theutterance \stepcounter{utterance}  
    & & & \multicolumn{2}{p{0.3\linewidth}}{
        \cellcolor[rgb]{0.9,0.9,0.9}{
            \makecell[{{p{\linewidth}}}]{
                \texttt{\tiny{[GM$|$GM]}}
                \texttt{Mode: passive\_examples} \\
\texttt{Accuracy: 0.080} \\
\texttt{Information requests: 3/15} \\
\texttt{Revealed passive examples: 4} \\
\texttt{Efficiency: 0.867} \\
\texttt{Slack: 12} \\
\texttt{Observed{-}pair consistency: 0.000} \\
\texttt{Parse errors: 6} \\
\texttt{Runtime errors: 0} \\
            }
        }
    }
    & & \\ \\

\end{supertabular}
}

\end{document}
